\documentclass[10pt,a4paper]{amsart}
\usepackage[margin=19mm]{geometry}
\usepackage{amsmath,amssymb,mathtools}
\usepackage[scr=rsfs,cal=euler]{mathalfa}
\usepackage{xfrac}
\usepackage[unicode,hidelinks,bookmarks=false]{hyperref}
\newcommand{\ie}{i.e.\ }
\newcommand{\cf}{cf.\ }
\newcommand{\resp}{respectively\ }
\newcommand{\Subsection}{\subsection}
\providecommand{\nobreakdash}{\nobreak\hbox{-}}
\setlength{\emergencystretch}{2em}
\multlinegap0pt
\usepackage{amssymb}
\usepackage{mathtools}
\usepackage{tikz}
\usepackage[shortlabels]{enumitem}
\usepackage{xcolor}
\usepackage{dsfont}
\usepackage[normalem]{ulem}
\newtheorem{lemma}{Lemma}
\newcommand{\scL}{\mathscr{L}}
\title{On cycles of lengths $8$ and $10$}
\author{David Conlon, Eion Mulrenin, Cosmin Pohoata}
\date{Source: arXiv:2603.24515v1 (CC BY 4.0)}
\begin{document}
\maketitle

\noindent\emph{Source and licence.} On cycles of lengths $8$ and $10$. Version arXiv:2603.24515v1 (CC BY 4.0), distributed by arXiv under the Creative Commons Attribution 4.0 International licence, http://creativecommons.org/licenses/by/4.0/, as recorded in the arXiv licence metadata for this exact version and shown on the abstract page https://arxiv.org/abs/2603.24515v1. Full licence notice: \texttt{licenses/arxiv-2603.24515-CC-BY-4.0.txt}.\par\smallskip

\noindent\textit{Editorial scope.} Counted unit: the article's lemma no-rectangles (source0.tex lines 367--370). The article's complete original proof of it is delivered with it (source0.tex lines 372--379, one \texttt{proof} environment of 8 lines). No mathematics is written, repaired or re-derived: the statement and the proof are the article's own lines, in the article's own notation, and no retained line is substituted or re-flowed.  No further source-local context is needed: every cross-reference of every retained line resolves inside the two delivered spans.  Nothing was omitted from any retained span, so the package carries no omission record.  No retained line contains a citation, so the wrapper carries no bibliography and no bibliography entry.  No retained line contains a figure environment, an image-inclusion command or drawing code, so no image is delivered and the package declares no asset dependency.  Editorial formatting only, in addition: an AMS \texttt{amsart} preamble that restates the article's own declarations and macros used by the retained lines, namely the environments \texttt{lemma} on separate counters, together with the standard packages those lines need.  The retained environments print in this document as Lemma 1 in order of appearance, whereas the article prints the same items with the numbering of its own full document.  The complete native source of the article as distributed in the arXiv e-print for this exact version is bundled unchanged as \texttt{sources/197153/source0.tex}.
    \begin{lemma}
    \label{lem: no-rectangles}
        If $H$ is $C_8$-free, then every auxiliary graph $G_\Pi^{i,j}(H)$ is $C_4$-free.
    \end{lemma}

    \begin{proof}
        A $4$-cycle in $G_\Pi^{i,j}(H)$ consists of two lines in $\scL_i(\Pi)$ and two lines in $\scL_j(\Pi)$ with all
        four cross-intersections active. 
        These four lines then form a copy of $C_8$ in $H$, as the four intersection points
        are distinct because lines with the same direction are parallel and distinct. 
        Therefore, a $C_4$ in some
        $G_\Pi^{i,j}(H)$ would produce a $C_8$ in $H$.
    \end{proof}

\end{document}
